\documentclass[aip,reprint,groupedaddress,onecolumn,superscriptaddress]{revtex4-2}
\usepackage{graphicx}
\usepackage{dcolumn}
\usepackage{bm}
\usepackage{amsmath}
\usepackage{siunitx}
\usepackage{booktabs}

\begin{document}

\title{Bulletpoint List for Main Text}
\date{\today}

\maketitle

\section{Experimental Framework: Co-Designed Quantum-Traceable Metrology}
\begin{itemize}
    \item \textbf{Dual-Wavelength Interferometric System}: We implemented a modified NIST LM10 Michelson interferometer operating simultaneously at \SI{780}{\nano\meter} (reference) and \SI{1762}{\nano\meter} (probe), with \SI{4.2}{\meter} optical path difference and air-floated mirror stage for vibration isolation.
    
    \item \textbf{Quantum Frequency References}: The \SI{1762}{\nano\meter} probe laser was stabilized to a ULE cavity referenced to a single trapped $^{138}$Ba$^+$ ion ($\delta f < \SI{1}{\kilo\hertz}$), while the \SI{780}{\nano\meter} reference used $^{87}$Rb saturation absorption spectroscopy ($\delta f \sim \SI{100}{\kilo\hertz}$).
    
    \item \textbf{Environmental Monitoring \& Data Collection}: Co-located BME280 sensors measured temperature, humidity, and pressure with \SI{1}{\hertz} sampling, collecting \num{145784} synchronized measurements from January to August 2025 across environmental ranges of \SIrange{15}{35}{\degreeCelsius} temperature, \SIrange{20}{45}{\percent} humidity, and \SIrange{965}{1000}{\hecto\pascal} pressure.
    
    \item \textbf{Statistical Methodology}: We employed ordinary least squares with HAC standard errors, validated through \num{2000}-sample block bootstrap and Prais-Winsten GLS for autocorrelation correction, using the model $n = \alpha_0 + \alpha_T T + \alpha_H H + \alpha_P P$.
\end{itemize}

\section{Key Results: Precision Coefficients \& Enhanced Sensitivity}
\begin{itemize}
    \item \textbf{High-Precision Refractive Index Coefficients}: We determined coefficients with part-per-billion precision: $\alpha_T = \SI{-8.8474e-7}{\per\degreeCelsius}$, $\alpha_H = \SI{-1.3152e-8}{\per\percent}$, $\alpha_P = \SI{+2.5950e-7}{\per\hecto\pascal}$, with $R^2 = 0.997$ and residual precision $\sigma_n = \num{1.82e-7}$.
    
    \item \textbf{Enhanced Humidity Sensitivity at 1762 nm}: Experimental measurement revealed \SI{12.4}{\percent} enhanced humidity sensitivity ($\partial n/\partial H = \SI{-1.3152e-8}{\per\percent}$) compared to standard models, consistent with proximity to H$_2$O absorption lines at \SI{1761.0405}{\nano\meter} and \SI{1762.4852}{\nano\meter} and Kramers-Kronig predictions.
    
    \item \textbf{Model Validation Performance}: Our empirical model showed superior performance with residual \num{-2.802e-14}$\pm$\num{1.837e-07} ($R^2 = 0.996$) versus Mathar model's \num{1.166e-06}$\pm$\num{1.930e-07} ($R^2 = 0.829$) on independent validation data.
    
    \item \textbf{Quantum Traceability Achievement}: The quantum reference contributed $\delta n < 10^{-12}$ uncertainty, establishing direct SI traceability from ambient refractive index to the single trapped $^{138}$Ba$^+$ ion frequency standard.
\end{itemize}

\section{Uncertainty Analysis \& System Performance}
\begin{itemize}
    \item \textbf{Comprehensive Uncertainty Budget}: Environmental sensor calibration (BME280) was identified as the dominant error source ($\delta n \sim 2.12\times10^{-7}$), with temperature calibration contributing $1.42 \times 10^{-10}$, humidity calibration $1.25 \times 10^{-10}$, and pressure calibration $9.11 \times 10^{-13}$.
    
    \item \textbf{Precision Metrics}: Total expanded uncertainty ($k=2$) was $4.24 \times 10^{-7}$, with residual precision (fit scatter) of $1.82 \times 10^{-7}$ and systematic uncertainty (sensor-limited) of $2.12 \times 10^{-7}$.
    
    \item \textbf{Statistical Robustness}: Analysis accounted for serial correlation through HAC standard errors, with effective sample size $N_{\text{eff}} \approx \num{39100}$ and Durbin-Watson statistic improved from 1.197 to 2.226 with Prais-Winsten correction.
    
    \item \textbf{System Architecture Validation}: The integrated quantum-traceable metrology system successfully bridged laboratory quantum references (GPS-disciplined oscillator, $^{138}$Ba$^+$ ion, $^{87}$Rb reference) with field-deployable sensor nodes, demonstrating part-per-billion precision ($\delta n \sim 1 \times 10^{-7}$) under realistic outdoor conditions.
\end{itemize}

\newpage

\section*{Main Text Structure Rationale}

This bulletpoint list organizes the main text content into three logical sections that mirror the experimental process:

1. \textbf{Experimental Framework}: Describes the "how" - the apparatus, methodology, and data collection approach that enables quantum traceability.

2. \textbf{Key Results}: Presents the "what" - the precise measurements, coefficients, and validation outcomes that constitute the core scientific findings.

3. \textbf{Uncertainty Analysis \& System Performance}: Provides the "how good" - the error analysis, precision metrics, and system validation that establish confidence in the results.

This structure follows the natural progression from experimental design to results to analysis, creating a clear narrative flow that supports the introduction's promise of bridging the "traceability gap" and sets the stage for the conclusion's discussion of impact and future directions.

\end{document}
